% answers_short.tex —— 2009 年·普通高等学校招生考试·江苏卷 · 简短答案速查
% 由 tools/make_short_answers.py 自 answers.tex 的 \gksolhead/\gkanswer 提取，
% 并与 2009/2009_数学真题_江苏_答案要点.md 逐题比对后生成，勿手改（改源 MD 后重跑）。

\begin{gknotice}
\textbf{使用说明}：本节只给最终答案便于快速核对，完整推导见解析版。答案为编写者独立推导，\textbf{非官方标准答案}。
\end{gknotice}

\gkpart{填空题}
\begin{itemize}[leftmargin=2.6em]
  \item[\textbf{1.}] $-20$。
  \item[\textbf{2.}] $3$。
  \item[\textbf{3.}] $(-1,11)$。
  \item[\textbf{4.}] $3$。
  \item[\textbf{5.}] $\dfrac15$。
  \item[\textbf{6.}] $\dfrac25$。
  \item[\textbf{7.}] $22$。
  \item[\textbf{8.}] $1:8$。
  \item[\textbf{9.}] $(-2,15)$。
  \item[\textbf{10.}] $m<n$。
  \item[\textbf{11.}] $4$。
  \item[\textbf{12.}] （1）（2）。
  \item[\textbf{13.}] $2\sqrt7-5$。
  \item[\textbf{14.}] $-9$。
\end{itemize}

\gkpart{解答题}
\begin{itemize}[leftmargin=2.6em]
  \item[\textbf{15.}] （1）$2$；（2）$4\sqrt2$；（3）证明见解析。
  \item[\textbf{16.}] （1）（2）证明见解析。
  \item[\textbf{17.}] （1）$a_n=2n-7$，$S_n=n^2-6n$；（2）$m=2$。
  \item[\textbf{18.}] （1）$y=0$ 或 $7x+24y-28=0$；（2）$\left(-\dfrac32,\dfrac{13}2\right)$、$\left(\dfrac52,-\dfrac12\right)$。
  \item[\textbf{19.}] U（原页未印 $m_A$、$m_B$ 的取值范围，按字面 (2) 无最大值、(3) 为「能」；若补入官方解答所用 $3\leqslant m_A\leqslant12$、$5\leqslant m_B\leqslant20$，则 (2) $m_A=12$、$m_B=20$、$h_0=\dfrac{\sqrt{10}}5$，(3) 为「不能」。详见解析与 ISSUE-1498。）
  \item[\textbf{20.}] （1）$(-\infty,-1]$；（2）$a\geqslant0$ 时最小值为 $-2a^2$，$a<0$ 时最小值为 $\dfrac{2a^2}3$；（3）分段解集见解析。
\end{itemize}

\gkpart{附加题}
\begin{itemize}[leftmargin=2.6em]
  \item[\textbf{21.}] A 证明见解析；B $\begin{bmatrix}-1&2\\2&-3\end{bmatrix}$；C $y=3x^2+6$（$x\in\mathbf R$）；D 证明见解析。
  \item[\textbf{22.}] （1）$y^2=2x$；（2）$2x+2y-1=0$；（3）$f(m)=\dfrac32\sqrt{m^2+4m}$（$m>0$）。
  \item[\textbf{23.}] （1）$T_{n^2}=\dfrac{n(n+1)(2n+1)}6+n^2(n^2-n)$，$P_{n^2}=1-\dfrac1n+\dfrac{2n^2+3n+1}{6n^3}$；（2）证明见解析。
\end{itemize}
